\newpage{\ }
\thispagestyle{empty}
\chapter*{Resumen}

\noindent El presente Trabajo Fin de Grado aborda el diseño, desarrollo e implementación de \textbf{UniHub}, un sistema informático distribuido e integral para la centralización, almacenamiento, consulta y gestión de la oferta de educación superior en España (universidades públicas y privadas, titulaciones de Grado, Máster y Doctorado, y planes de estudio desglosados del Boletín Oficial del Estado - BOE).

El proyecto se estructura en cuatro fases de ingeniería perfectamente delimitadas e interconectadas:
\begin{enumerate}
    \item \textbf{Fase 1 (Rastreador / Crawler):} Desarrollado en Python, obtiene el catálogo oficial, filtra titulaciones vigentes y analiza las resoluciones del BOE. El análisis de PDF construye un modelo de páginas, posiciones y tablas para delimitar cada titulación y extraer su plan curricular. La ejecución registra su trazabilidad, respeta \texttt{robots.txt} en las fuentes web y admite una planificación Cron configurable mediante \texttt{CRAWLER\_CRON\_SCHEDULE}.
    \item \textbf{Fase 2 (Base de Datos PostgreSQL y API REST):} Modelo relacional DDL con rol de solo lectura de API (\texttt{unihub\_api\_user}), pipeline de carga ETL, ordenación prioritaria (públicas primero, luego privadas), soporte de métricas físicas individuales de contenedores cgroup y servicio REST en FastAPI.
    \item \textbf{Fase 3 (Portal Web Frontend):} Aplicación Web SPA moderna creada con React y Vite bajo la identidad visual de la Universidad de Cádiz (UCA). Incluye paginación configurable (5, 10, 20, 50 y 100 por página), ocultación de códigos internos, geolocalización por fórmula de Haversine con distancias integradas en tarjetas, sección de universidades destacadas (top 6 más visitadas), apartado "Sobre Nosotros" (TFG Alejandro Ramos Rodríguez, UCA) y panel de administración aislado en la ruta manual \texttt{/admin}.
    \item \textbf{Fase 4 (Contenerización y Orquestación Docker):} Sistema de 4 contenedores aislados (\texttt{unihub\_db}, \texttt{unihub\_crawler}, \texttt{unihub\_api}, \texttt{unihub\_www}) orquestados con Docker Compose con garantía de persistencia permanente de datos tras apagar o reiniciar los servicios.
\end{enumerate}

\vspace{.5cm}

\noindent \textbf{Palabras clave:} UniHub, Universidad de Cádiz (UCA), Rastreador Python, BOE, Doctorado, PostgreSQL, FastAPI, React, Geolocalización Haversine, Docker Compose.
